% Transcription — fonds Alexandre Grothendieck, shelfmark 2, batch 2 (pages 21-40).
% Established from the facsimile archives/batches/2/batch-02.pdf (a mirror of
% https://grothendieck.umontpellier.fr/2.pdf), per the skill
% transcribe-grothendieck. The batch file carries no cover sheet and opens at
% the archivists' page 21: PDF page N is archivists' page N+20. Checked against
% the pencilled numbers bottom-left ("21", "22", "24", "27", "28", "32", "34",
% "36", "40" read on the corresponding leaves).
%
% Pass: Opus 5 (claude-opus-5), 2026-09-19 — first pass, unchecked against the
% pages by a human.
%
% What the batch is. Two distinct runs, on two kinds of paper, with covers in
% his hand between them.
%
%   Pages 21-24, ink on white leaves, fast hand. The abstract theory: a
%     lambda-ring K with augmentation eps, its gamma-filtration K^(i), the ring
%     CK generated by Chern classes, the associated graded GK, and the
%     comparison map between them. Page 21 closes the proof that phi_B psi_B is
%     multiplication by (-1)^{i-1}(i-1)!, then asks whether phi_K is bijective
%     and answers no, on the example K = Z[L]/(L-1)^n, which he notes is
%     K(P^{n-1}). Page 22 is the free lambda-ring B on generators X_s, Y_t, its
%     weight filtration B^(k), and the lemma that the total class c~ is
%     additive and multiplicative. Page 24 is the proposition that lambda_t(x)
%     polynomial of degree n forces gamma^i(x - eps(x)) = 0 for i > n, with a
%     remark and a corollary.
%   Pages 27-29, ink on blue-grey leaves. The axiomatic set-up restated: K
%     Z-augmented with a family (E_s) of elements of finite rank n_s with
%     invertible top exterior power, the determinant map K -> K*, the Chern
%     homomorphism c : K -> A~, then on page 28 the ring C*K by generators
%     c^i(x) and four families of relations, and phi_K : CK -> GK sending
%     c^i(x) to gamma^i(x - eps(x)) mod K^(i+1). Page 29 is three lines and
%     stops.
%   Pages 32-37, pencil then ink on typescript versos, under the cover of page
%     31. Geometry: D_parf(f) and K_parf(f), the same with supports in a closed
%     Z or a family Phi, the limit over a system of infinitesimal
%     neighbourhoods, then K_Phi(X) and A_Phi(X) with Chern classes c^i_Phi
%     between them, their covariance, and on page 36 the sentence that R.R. is
%     what interprets the discrepancy.
%   Page 40, under the cover of page 39: the Grothendieck-Riemann-Roch square
%     for a proper f : X -> Y, the formula (1) and its variant (1 bis), the
%     refinement (2) for a closed immersion, and the example of a fibration.
%
% Skipped, without description, as the skill requires: pages 23, 25, 33, 35 and
% 38 (typescripts that are not his and do not carry mathematics in his hand —
% 33 and 35 carry only line-by-line corrections to the typed text); pages 26
% and 30 (blank leaves). Pages 31 and 39 are otherwise blank sheets inscribed
% in his hand naming the runs that follow: they get no \page, and their words
% are the section headings marked as his.
%
% Notation fixed once here and held. The lambda-ring is K, its augmentation
% \varepsilon, its operations \lambda^{i} and \gamma^{i}; the gamma-filtration
% is K^{(i)} and B^{(k)} on the free ring B. The ring generated by the Chern
% classes is CK (C^{*}K where he writes the star on page 28), the associated
% graded GK — he also writes \mathrm{Gr}\,K on page 29 and both are kept as
% written. The comparison maps are \varphi_{B}, \psi_{B}, \varphi_{K}. His
% capital A with a tilde is \widetilde{A}, the total class \widetilde{c}, the
% graded pieces A^{(i)}. In the geometric run the derived categories are
% D_{\mathrm{parf}} (the subscript is read "parf" and is legible as such on
% page 32), the families of supports \Phi, the Chow-type groups A_{\Phi}(X),
% the Chern character \operatorname{ch} and the Todd class
% \operatorname{Todd}; on page 40 the two vertical composites are
% \tau_{h}(X), \tau_{h}(Y) and the relative tangent object is T_{Y/X}.
%
% Readings flagged and not settled. On page 21 the first symbol of the
% demonstration, read CB; the marginal note running up the left edge, of which
% only the ideal and the ring Z[\xi]/\xi^{n} come through; the numbering of the
% proposition and of the corollary. On page 22 the word before "poids" in the
% definition of B^{(k)}, and the whole of the last three lines, where a
% universal polynomial in the \gamma^{k}(x), \gamma^{h}(y) is announced and
% the argument breaks off. On page 24 the two struck blocks and the words
% between "on peut supposer" and the ring Z[\lambda^{i}(X)]. On page 27 the
% word underlined after "Une deuxième", read "explicite", and the last word of
% the cover of page 31, read "supports". On page 28 the long N.B. at the foot,
% which the pen changes for and which is lost below the third word. On page 32
% the object of which R i_{*} is taken, and the qualifier of the neighbourhoods
% Z_n. On page 34 the parenthesis after A_{\Phi}(X), read "pseudo-anneaux",
% and the three numbered cases at the foot. On page 36 almost all of the
% connective prose. On page 37 the four labels on the vertical arrows, all of
% the form "isom. si X'/X rig.", and the scribbled-out block at mid-page. On
% page 40 the words after "f une injection" and after "f une surjection", and
% the closing five lines, which the ink of the verso crosses.
\documentclass[11pt,a4paper]{article}
% Le préambule commun à toutes les transcriptions du fonds Grothendieck.
%
% Il tient en une idée : **l'incertitude se compose, elle ne se résout pas.**
% Une transcription qui lisse un mot douteux a détruit la seule chose pour
% laquelle on la faisait — un lecteur doit pouvoir savoir, à chaque endroit,
% ce qui était sur le feuillet et ce qui a été ajouté. Les six macros qui
% suivent sont donc l'appareil critique, et non de la décoration.
%
% Les mêmes commandes sont reconnues par `scripts/render.mjs`, qui produit la
% vue de lecture du site. Une macro ajoutée ici doit l'être aussi là-bas,
% faute de quoi le rendu échouera bruyamment — ce qui est voulu : un rendu
% silencieusement faux est pire qu'un rendu absent.

\ProvidesPackage{grothendieck}[2026/08/08 Transcriptions du fonds Grothendieck]

\usepackage{amsmath,amssymb,amsthm}
\usepackage{mathrsfs}
\usepackage{stmaryrd}
% Les diagrammes commutatifs sont partout dans ce fonds : c'est la langue
% même de la géométrie algébrique de Grothendieck, et une transcription qui
% les rendrait en prose ne transcrirait rien.
\usepackage{tikz-cd}
% Français par défaut — c'est la langue des feuillets — mais l'anglais est
% chargé aussi : la traduction et le résumé sont des documents anglais, et la
% typographie française (espaces avant les deux-points) y serait une faute.
% Ces documents basculent par \selectlanguage{english} en tête de corps.
\usepackage[english,french]{babel}
\usepackage{fontspec}
\usepackage{geometry}
\geometry{a4paper,margin=25mm,marginparwidth=28mm}
\usepackage{marginnote}
\usepackage{xcolor}
% ulem, not soul. `soul`'s \st and \ul are built on a character-by-character
% scan that XeLaTeX's font machinery breaks: the document fails with an
% undefined control sequence rather than degrading. ulem does the same two jobs
% and is Unicode-engine safe. `normalem` stops it hijacking \emph, which the
% transcriptions use for Grothendieck's own emphasis.
\usepackage[normalem]{ulem}
\usepackage{hyperref}
\hypersetup{colorlinks=true,linkcolor=black,urlcolor=black,pdfborder={0 0 0}}

% --- Métadonnées du lot -----------------------------------------------------
% Reprises telles quelles par le script de rendu, qui les affiche en tête de
% la vue de lecture. Elles fixent ce que le fichier prétend couvrir : sans
% elles, une transcription flotte sans point d'attache dans un fonds de seize
% mille feuillets.

\newcommand{\folder}[1]{\gdef\grothFolder{#1}}
\newcommand{\batch}[1]{\gdef\grothBatch{#1}}
\newcommand{\leaves}[2]{\gdef\grothFirst{#1}\gdef\grothLast{#2}}
\newcommand{\foldertitle}[1]{\gdef\grothFoldertitle{#1}}
\gdef\grothFolder{?}\gdef\grothBatch{?}\gdef\grothFirst{?}\gdef\grothLast{?}\gdef\grothFoldertitle{}

% --- Le résumé --------------------------------------------------------------
%
% Il ouvre la lecture modernisée, sous le titre « Résumé », et il est composé
% en petit. Ce n'est pas un ornement : c'est le seul passage du document qui
% ne soit pas des mathématiques, et un lecteur doit pouvoir le reconnaître
% avant de l'avoir lu — puis le sauter s'il connaît déjà le sujet, ou s'y
% arrêter s'il ne le connaît pas. Le filet qui le suit marque la reprise du
% corps.

\newenvironment{resume}
  {\small\setlength{\parskip}{0.45em}}
  {\par\medskip\hrule\medskip}

% --- Mots-clés ---------------------------------------------------------------
%
% En anglais, en clôture du résumé de la lecture modernisée : le vocabulaire
% moderne sous lequel chercher ce que les feuillets construisent. Le manifeste
% du site (`npm run manifest`) les extrait pour étiqueter la cote entière —
% cette ligne est la source unique des tags, il n'y a pas de fichier à côté.
\newcommand{\keywords}[1]{\par\smallskip\noindent
  {\footnotesize\textsc{Keywords} --- \textit{#1}}\par}

% --- L'appareil critique ----------------------------------------------------

% Le numéro de feuillet, en marge. C'est l'ancre qui permet de régler un
% désaccord sur une formule en regardant *un* feuillet plutôt qu'une chemise
% de six cents. Le site s'en sert aussi pour tourner les pages du fac-similé
% quand on fait défiler la transcription.
\newcommand{\leaf}[1]{%
  \marginnote{\footnotesize\color{gray!70}\textsf{#1}}%
  \label{leaf:#1}\ignorespaces}

% Illisible : on ne devine pas. Un mot inventé qui se lit comme les autres est
% le pire résultat possible.
\newcommand{\ill}{\textcolor{red!60!black}{[\,\ldots\,]}}

% Lecture douteuse : le mot est donné, et signalé comme incertain.
\newcommand{\uncertain}[1]{\uline{#1}}

% Ajout de l'éditeur — une lettre manquante, un mot sous-entendu.
\newcommand{\add}[1]{\textcolor{blue!50!black}{[#1]}}

% Biffé par Grothendieck lui-même. Ce qu'il a rayé fait partie du document :
% une rature dit souvent où la pensée a changé de direction.
\newcommand{\struck}[1]{\sout{#1}}

% Note du transcripteur, sur l'état matériel du feuillet.
\newcommand{\note}[1]{\par\smallskip{\small\color{gray!60!black}\textit{[#1]}}\par\smallskip}

% Note marginale de Grothendieck, distincte du corps du texte.
\newcommand{\marginal}[1]{\par\smallskip\noindent\hspace{1em}%
  {\small\color{gray!60!black}#1}\par\smallskip}

% --- Titre ------------------------------------------------------------------

\AtBeginDocument{%
  \begin{center}
    {\large\bfseries Fonds Alexandre Grothendieck --- cote n\textsuperscript{o}~\grothFolder}\\[2pt]
    {\itshape\grothFoldertitle}\\[4pt]
    {\small feuillets \grothFirst--\grothLast\ (lot \grothBatch)}\\[2pt]
    {\footnotesize\color{gray!60!black}Transcription établie d'après le fac-similé numérisé
     par l'Université de Montpellier.\\
     Les lectures incertaines sont soulignées, les illisibles notées [\,\ldots\,],
     les ajouts entre crochets.}
  \end{center}
  \vspace{1em}}

\endinput


\folder{2}
\batch{2}
\pages{21}{40}
\dating{[1957]}
\watermark{Édition de démonstration}
\foldertitle{Théorème de Riemann-Roch / RR [Riemann-Roch] Notes Antiques (antérieur[es à] 1958) : notes manuscrites (s.d.).}

\begin{document}

\section*{Comparaison de $CK$ et de $GK$}

\page{21}

\emph{Démonstration.} Comme \uncertain{$CB$} $\to CK$ et $GB \to GK$ sont
surjectifs, on est ramené au cas de $B$.

Soit $\xi^{i} \in C^{i}(B)$ provenant de $P \in B^{(i)}$, on a
\[
  c_{B}^{i}(\xi^{i}) \;=\; C^{i}_{B}(P) ,
\]
\ill{} d'où
\[
  \varphi_{B}\psi_{B}(\xi^{i}) \;=\; \varphi_{B}\bigl(C^{i}_{B}(P)\bigr)
  \;=\; c^{i}_{B}(P) \;\equiv\; \struck{c^{i}_{B}}\ \struck{dans}\
  \gamma^{i}\bigl(P - \varepsilon(P)\bigr) \pmod{B^{(i+1)}} ,
\]
et
\[
  \gamma^{i}\bigl(P - \varepsilon(P)\bigr) \;-\; (-1)^{i-1}(i-1)!\,P
  \;\in\; B^{(i+1)} ,
\]
soit \ill{}
\[
  \varphi_{B}\psi_{B} \;=\; (-1)^{i-1}(i-1)! ,
\]
cqfd.

\uncertain{Proposition 3.} \ill{} $\varphi_{B}$ \ill{} bijectif.

\emph{Corollaire.} $\varphi_{K}$ \ill{} toujours surjectif \ill{}

\emph{Question.} Est-ce que $\varphi_{K}$ est toujours bijectif ??? Non
\note{la réponse est encadrée sur le feuillet, avec le renvoi à l'exemple qui suit}

\emph{Exemple.} Soit $K = \mathbb{Z}[L]/(L-1)^{n}$, avec
$\lambda_{t}(L) = 1 + Lt$ \ill{} , donc
\[
  CK \;=\; \mathbb{Z}[\xi] \big/ \bigl(\text{idéal engendré par } \ill{}\bigr),
  \qquad
  (1+\xi)^{\ill{}} \;=\; 1 + (-1)^{n-1}(n-1)!\,\xi^{n} + \cdots
\]
\marginal{c'est-à-dire $K = K(P^{\,n-1})$}
Cependant, \ill{} la filtration sur $K$ est la filtration \ill{} par les
puissances de $(L-1)$, donc \ill{} sans torsion, tandis que $CK$ a de la
torsion. \ill{}

\marginal{Mais si \ill{} l'idéal était \ill{} : $P(\ill{})$ \ill{}
$CK = \mathbb{Z}[\xi]/\xi^{n}$ \ill{}}
\note{la note court en diagonale le long du bord gauche du feuillet ; seuls
l'idéal et l'anneau $\mathbb{Z}[\xi]/\xi^{n}$ s'en laissent lire}

\emph{Proposition} \add{(}à prouver directement\add{?)}. Soit $K$ un
$\lambda$-anneau, $(X_{s})_{s \in S}$ une famille d'éléments \ill{} ;
conditions équivalentes :
\begin{enumerate}
  \item \ill{} engendré par les $X_{s}$, \ill{} un déterminant \ill{} ;
  \item $c^{i}(X_{s}) = 0$ si $i > \varepsilon(X_{s})$ \ill{}
\end{enumerate}
\note{tout le bloc est barré d'une grande croix ; il est repris, sous une autre
forme, à la page 27}

\page{22}

\[
  B \;=\; \mathbb{Z}\bigl[\,(\lambda^{i}(X_{s}))_{s \in S,\; 1 \leqslant i \leqslant \varepsilon(s)},\;
  (\lambda^{d}(Y_{t}))_{t \in T,\; d \geqslant 1}\,\bigr] ,
\]
\[
  CB \;=\; \mathbb{Z}\bigl[\,(\xi^{i}_{s})_{s \in S,\; 1 \leqslant i \leqslant \varepsilon(s)},\;
  (\eta^{d}_{t})_{t \in T,\; d \geqslant 1}\,\bigr] .
\]

Filtration sur $B$ : \struck{pour} $B^{(k)} =$ sous-anneau \ill{} engendré par
les \uncertain{monômes} $P$ \ill{} de \uncertain{poids} $\geqslant k$ en les
\[
  \gamma^{i}\bigl(X_{s} - \varepsilon(X_{s})\bigr), \qquad
  \gamma^{d}\bigl(Y_{t} - \varepsilon(Y_{t})\bigr) ,
\]
\ill{} convention : \ill{} définis par les $\varepsilon(s)$, $\varepsilon(t)$.

\struck{\ill{}}

\emph{Lemme.} Posons $c^{i}(x) = \gamma^{i}\bigl(x - \varepsilon(x)\bigr)
\bmod B^{(i+1)}$. Alors
\[
  x \;\longmapsto\; \widetilde{c}(x) \;=\;
  \bigl[\,\varepsilon(x),\; 1 + c^{1}(x)\,t + \cdots \,\bigr]
\]
\ill{} On a
\[
  \widetilde{c}(x+y) \;=\; \widetilde{c}(x)\,\widetilde{c}(y)
\]
\ill{} dans le $\lambda$-anneau \uncertain{$\widetilde{GB}$} \ill{} : posant
$x' = x - \varepsilon(x)$, $y' = y - \varepsilon(y)$, on a
$\gamma_{t}(x' + y') = \gamma_{t}(x')\,\gamma_{t}(y')$.

\[
  \widetilde{c}(xy) \;=\; \widetilde{c}(x) \ast \widetilde{c}(y)
\]
\ill{} : a) $x = 1$ ou $y = 1$, trivial ; b) $x$ et $y$ \ill{} Alors \ill{}

\struck{$c^{i}(xy) = Q_{i}(\ill{})$}
\[
  \widetilde{c}(xy) \;=\; Q_{i}\bigl(c^{1}(x), \ldots, c^{i}(x);\;
  c^{1}(y), \ldots, c^{i}(y)\bigr) \qquad \text{soit} \ \ill{}
\]
car \ill{}
\[
  \gamma^{i}(xy) \;=\; Q_{i}\bigl(\gamma^{1}(x), \ldots, \gamma^{i}(x);\;
  \gamma^{1}(y), \ldots, \gamma^{i}(y)\bigr) .
\]
Mais $\gamma^{i}(xy) =$ polynôme universel en \ill{} les $\gamma^{k}(x)$,
$\gamma^{h}(y)$, \ill{} et \ill{} soit \ill{} $\geqslant i$, et
\struck{\ill{}} $Q_{i}$
\note{la page s'arrête ici : le polynôme universel est annoncé et sa forme
n'est pas écrite}

\page{24}

\emph{Proposition.} Soit $x \in \widetilde{A}$ tel que $\lambda_{t}(x)$ soit un
polynôme de degré $\leqslant n$, \ill{} $\varepsilon(x) = n$. Alors
$\gamma^{i}(x) = 0$ pour $i > n$, i.e. \ill{}

En effet, $\lambda^{n}(x)$ \ill{} et le fait que \ill{} $x$ \ill{} nilpotent
\ill{}
\[
  \gamma^{i}\bigl(x - \varepsilon(x)\bigr) \;=\; 0 \qquad \text{pour } i > n ,
\]
alors \ill{}
\[
  (i-1)!\;x^{i} \;=\; 0 \qquad \text{pour } i > n ,
\]
pour $[\,1,\; 1 + x^{1}\,\ill{}\,]$ dans \ill{} arithmétique \ill{} de
$\widetilde{A}$, donc \ill{} $> n$. \ill{}

\emph{Remarque.} Posons, pour
$x = [\,n,\; 1 + \ill{} + x^{n} + x^{n+1}\,]$, \ill{} soit
$n!\;x^{n+1} = 0$ \ill{} ; il faut et il suffit que \ill{} $n!$ \ill{} Posant
$Y = [\,n+1,\; 1 + X^{n+1}\,]$, d'où $x = X + Y - (n+1)$, qui \ill{}
\[
  D^{i}_{t}\bigl(\lambda_{t}(X)\,\lambda_{t}(Y)\bigr)\big|_{t=1} \;=\; 0
  \qquad \text{pour } 0 \leqslant i \leqslant m .
\]

\emph{Corollaire.} Soit $x$ un élément \ill{} Soit $n = \varepsilon(x)$, et
$m \geqslant \ill{}$ \ill{} $\gamma^{i}(x)$ \ill{} Alors
$\gamma^{m}(x - \ill{})$ \ill{} $\lambda^{m}(X)$ ($i > \ill{}$) \ill{}
$m+1$ \ill{} les $\gamma^{i}(X - n)$.

\struck{\ill{}}

On peut supposer $A = \mathbb{Z}\bigl[(\lambda^{i}(X))_{i \geqslant \ill{}}\bigr]$,
\ill{} $\varepsilon(X) = m$.

\emph{Démonstration.} Soit $I$ l'idéal \ill{} par les $\lambda^{i}(X)$
($i > n$), \ill{} que $\gamma^{m}(X - m)$ \ill{} $I$ \ill{} les
$P\bigl(\gamma^{i}(X - m)\bigr)$, $P$ \ill{} de poids $m+1$. Sur
$B = \mathbb{Z}\bigl[(\lambda^{i}(X))_{1 \leqslant i \leqslant m}\bigr]$ :
\ill{} $\gamma^{m}(X - m)$ \ill{} de « filtration » $\geqslant m+1$ (pour
$m \geqslant n$).

\section*{Les axiomes : $K$ augmenté, les $E_{s}$, le déterminant et les classes de Chern}

\page{27}

$K$ un $\lambda$-anneau $\mathbb{Z}$-augmenté,
\[
  \varepsilon : K \longrightarrow \mathbb{Z} ,
\]
muni d'une famille d'éléments $(E_{s})_{s \in S}$ satisfaisant
\[
  \begin{cases}
    \varepsilon(E_{s}) = n_{s} \geqslant 0, \\[2pt]
    \lambda^{i} E_{s} = 0 & \text{si } i > n_{s}, \\[2pt]
    \lambda^{n_{s}} E_{s} \text{ inversible},
  \end{cases}
\]
\ill{} pour les $\lambda^{i}$ et (avec $\varepsilon = 1$) pour $\sigma$ :
$\lambda^{i}$ \ill{}

$[\,$\ill{} engendré \ill{} $K\,]$, \ill{} et \ill{}
\[
  K \;\xrightarrow{\ \text{dét}\ }\; K^{*}
\]
tel que
\[
  \text{dét}(E_{s}) \;=\; \lambda^{n_{s}}(E_{s}) .
\]
$[\,$ce qui \ill{} détermine \ill{}$\,]$

Une deuxième \struck{donnée} \uncertain{explicite} \ill{} des structures
précédentes \ill{} « dans » un $\widetilde{A}$ \ill{} , \ill{}
\struck{\ill{}} \ill{} de $\lambda$-\ill{} \ill{}
\[
  c : K \longrightarrow \widetilde{A}
\]
$[\,$\ill{} défini par $\lambda$, \ill{} commutatifs \ill{} sur les
$E_{s}\,]$ tel que
\[
  c^{i}(E_{s}) \;=\; 0 \quad \text{pour } i > n_{s} ,
  \qquad
  c^{1}(E) \;=\; c^{1}(\text{dét } E) .
\]
\ill{} comparaison ?? \ill{} $\lambda^{n_{s}}(E_{s})^{-1}$ \ill{} les
$\lambda^{i}$ de \ill{} $E_{s}$ \ill{} $A$ universel,

\marginal{En particulier, ceci \uncertain{définit} pour $n_{s} = 1$ :
$\widetilde{E}_{s}$, $\widetilde{c}(E_{s}) = [\,1,\; 1 + c^{1}(E_{s})\,]$ et
$c^{1}(E_{s})$ additive : $c_{1}(E_{s}E_{t}) = c_{1}(E_{s}) + c_{1}(E_{t})$}

\page{28}

\ill{} noté $C^{*}K$, défini par les générateurs
\[
  c^{i}(x) \qquad i \geqslant 1
\]
et les relations
\[
  \begin{cases}
    c^{i}(1) = 0 & i \geqslant 1, \\[3pt]
    c^{n}(x+y) = \displaystyle\sum_{i+j=n} c^{i}(x)\,c^{j}(y)
      & x, y \in K,\ n \geqslant 1, \\[8pt]
    c^{n}(x \ast y) = Q_{n}\bigl(c^{1}(x), \ldots, c^{n}(x);\;
      c^{1}(y), \ldots, c^{n}(y)\bigr)
      & x, y \in I(K),\ n \geqslant 1, \\[5pt]
    c^{n}(E_{s}) = 0 & n > n_{s},
  \end{cases}
\]
\ill{} pour tous les $E_{s}$.

Notons que l'on a
\[
  \gamma^{n}(E_{s} - n_{s}) \;=\; 0 \qquad \text{pour } n > n_{s} ,
\]
$[\,$\ill{} de $\lambda^{i}(E_{s}) = 0$ pour $i > n_{s}\,]$, donc \ill{}
les \ill{} \struck{\ill{}} satisfont aux conditions de \ill{} admissibles
\ill{} pour les $E_{s}$, \ill{}
\[
  \varphi_{K} : CK \longrightarrow GK, \qquad
  c^{i}(x) \;\rightsquigarrow\; \gamma^{i}\bigl(x - \varepsilon(x)\bigr)
  \bmod K^{(i+1)} ,
\]
évidemment surjectif.

$[\,$N.B. plus la famille $(E_{s})$ \ill{} , plus \ill{} $CK$ devient plus
petit, \ill{} plus \ill{} que $\varphi_{K}$ soit un isomorphisme \ill{}$\,]$

\marginal{À voir : le $\ast$ de l'itération \ill{} de
$P\bigl(\gamma^{i}(x - \varepsilon(x))\bigr)$ \ill{}
$P\bigl(\gamma^{i}(E_{s} - \varepsilon(E_{s}))\bigr)$ \ill{}
$\varepsilon(E_{s})$ \ill{} ??}

N.B. \ill{}
\note{la longue note du bas du feuillet, d'une plume plus fine, ne se laisse
lire qu'au premier mot}

\page{29}

D'ailleurs, \ill{} fois que tout homomorphisme \ill{} de Chern
\[
  \chi : K \longrightarrow \widetilde{A}
\]
$[\,$qui « respecte » les $E_{s}$ \ill{}$\,]$ transforme les éléments de
$K^{(i)}$ en éléments de $A^{(i)}$, d'où \struck{(additif)} \ill{}
\[
  \psi : \mathrm{Gr}\,K \longrightarrow \struck{A}\ G(\widetilde{A}) \simeq A\,!
\]
$\bigl[\simeq A$ additivement \ill{} $\neq 0\bigr]$

En particulier \ill{} on a \ill{} les
\[
  \psi : \mathrm{Gr}\,K \longrightarrow \ill{} ,
  \qquad
  K \longrightarrow CK\,!
\]
d'ailleurs \ill{}
\note{le feuillet s'arrête ici ; les deux tiers inférieurs sont vierges}

\section*{Construction des Chern covariants, et $K^{\bullet}$ et $A^{\bullet}$ à supports}

\note{titre pris du feuillet 31, qui ne porte que ces mots, de sa main :
« Construction des Chern covariants, et $K^{\bullet}$ et $A^{\bullet}$ à
\uncertain{supports} »}

\page{32}

Soit $f : X \to Y$ localement de type fini. On définit $D_{\mathrm{parf}}(f)$
comme la catégorie \struck{des} \uncertain{hégélienne} \ill{} complexes de
$\mathcal{O}_{X}$-Modules \struck{parfaits} parfaits relativement à $f$, et
$K_{\mathrm{parf}}(f)$ en conséquence.

Soit $Z \subset X$ un fermé \ill{} On définit $D_{\mathrm{parf},Z}(X)$ comme la
catégorie de \ill{} complexes de $\mathcal{O}_{X}$-Modules parfaits \ill{}
\uncertain{conjoint} \ill{} d'où $K_{Z}(X)$.

On définit de même, $\Phi$ \ill{} famille :
\[
  D_{\mathrm{parf},\Phi}(X), \quad K_{\Phi}(X), \qquad
  D_{\mathrm{parf},\Phi}(X) = \varinjlim_{Z \in \Phi} D_{\mathrm{parf},Z}(X),
  \quad
  K_{\Phi}(X) = \varinjlim_{Z \in \Phi} K_{Z}(X) .
\]

Soit $i : Z \to X$ \ill{} $f$ \ill{}
\[
  R\,i_{*} : D_{\mathrm{parf}}(\ill{}) \longrightarrow D_{Z}(X),
  \qquad
  i_{!} : K_{\ill{}}(\ill{}) \longrightarrow K'_{Z}(X) .
\]

Soit \struck{\ill{}} $i_{n} : Z_{n} \to X$ \ill{} \uncertain{voisinages}
\ill{} de $Z$ dans $X$, \ill{} alors
\[
  \varinjlim_{n} D_{\mathrm{parf}}(i_{n}) \longrightarrow D_{Z}(X) ,
\]
le foncteur \ill{} pleinement fidèle \ill{}

\page{34}

Définir $K_{\Phi}(X)$ ($X$ préschéma), définir \add{en} substituant
\struck{et} l'anneau de Chow
\[
  \struck{A}\ A_{\Phi}(X)
\]
(qui \ill{} être les \add{(}pseudo-\add{)}anneaux \ill{} par voie d'éliminer
$1$) et les classes de Chern \ill{} « l'anneau de Chow \emph{ordinaire}
whatever that means »
\note{les trois derniers mots sont en anglais sur le feuillet}
\[
  c^{i}_{\Phi} : K_{\Phi}(X) \longrightarrow A_{\Phi}(X) \qquad [\,i > 0\,]
\]
qui \ill{} Les conditions de multiplicativité \ill{} ordinaire \ill{}
\[
  c^{n}_{\Phi}(x+y) \;=\; c^{n}_{\Phi}(x) + c^{n}_{\Phi}(y)
  \;+\; \sum_{\substack{p+q=n \\ p,\,q > 0}} c^{p}_{\Phi}(x)\,c^{q}_{\Phi}(y)
  \qquad \bigl(x, y \in K_{\Phi}(X)\bigr)
\]
et \struck{\ill{}} être compatibles avec les homomorphismes canoniques

\begin{tikzcd}[column sep=large, row sep=small]
K_{\Phi}(X) \arrow[r] \arrow[d, "c_{\Phi}"'] & K(X) \arrow[d, "c"] \\
A_{\Phi}(X) \arrow[r] & A(X)
\end{tikzcd}

\ill{} la démonstration est \ill{} les multiplicités \ill{}
\struck{Intégrer} étudier \ill{} canoniques\add{)}. Cas :

\begin{enumerate}
  \item $X$ \struck{est} lisse \ill{} définition \ill{} de \ill{} , mais en
    prenant \ill{} les cycles sur $X$ à support $\subset \Phi$ ;
  \item $X$ général : on \ill{} manière de définir $A_{\Phi}(X)$ \ill{} par
    la $\lambda$-procédure \ill{} si une filtration sur \ill{}
    $K_{\Phi}(X)$ ?? ;
  \item Classes de Chern \uncertain{cohomologiques} ??
\end{enumerate}

\page{36}

On voit les voisinages \ill{} les plus favorables \ill{} les $K_{\Phi}(X)$,
$A_{\Phi}(X)$ \ill{} \struck{\ill{}} \ill{} (multiplications parfaites)
\ill{} , \ill{} $\to B$. \ill{} Utile\add{)} \ill{} alors, les
classifications, définies dans \ill{}
\[
  K^{\Phi}_{\ill{}}(X/S) \;=\; K^{\Phi}_{\ill{}}(f),
  \qquad
  A^{\Phi}_{\ill{}}(X/S) \;=\; A^{\Phi}_{\ill{}}(f),
\]
\ill{} « covariants » \ill{} \struck{\ill{}} \ill{} « covariants » \ill{}
classes de Chern
\[
  \operatorname{ch}^{X,\Phi}_{\ill{}} : K_{\ill{}}(X/S) \longrightarrow
  A_{\ill{}}(X,S),
  \qquad
  K^{\Phi}_{\ill{}}(X/S), \quad A^{\Phi}_{\ill{}}(X/S) \ \text{covariants} ,
\]
\ill{} les types de propriétés \ill{}
\[
  \operatorname{ch}^{Y/S}_{\ill{}}\bigl(g_{\ill{}}(x)\bigr)
  \;=\; g_{*}\bigl(\operatorname{ch}^{X/S}_{\ill{}}(x)\bigr) ,
\]
\struck{$g_{\ill{}}\bigl(\operatorname{ch}^{X}_{\ill{}}(x)\bigr)
= \operatorname{ch}^{\ill{}}_{\ill{}}(\ill{})$}
\ill{} plus avec l'image directe. C'est-à-dire compatible dans \ill{}

Mais « R.R. » s'interprète comme \ill{} donnant \ill{} dans certains cas
\ill{} fait calculé \ill{}
\[
  \bigl[\,X/S \ \text{parfait} \ \ill{}\,\bigr]
\]
\struck{d'image} $\operatorname{ch}^{X/S}_{\ill{}}$ \ill{} comme \ill{}
à \ill{} \add{(}comme que\add{)} explicitables \ill{}
\[
  K^{\bullet}(X) \longrightarrow K_{\ill{}}(X/\Sigma) ,
\]
\ill{} utilisant \ill{} $K^{\bullet}(X)$ \ill{}

\page{37}

\begin{tikzcd}[column sep=small, row sep=small, nodes={font=\scriptsize}]
K^{\bullet}(i_{n}) \arrow[r, "u"] \arrow[d, "\alpha"'] & \varinjlim_{n} K^{\bullet}(i_{n}) \arrow[r] \arrow[d, "\beta"] & K^{\bullet}_{X}(X') \arrow[r] \arrow[d, "\gamma"] & K^{\bullet}(X') \arrow[d, "\delta"] \\
K_{\bullet}(X_{n}) \arrow[r, "u'"'] & \varinjlim_{n} K_{\bullet}(X_{n}) \arrow[r] & K^{X}_{\bullet}(X') \arrow[r, "w'"'] & K_{\bullet}(X')
\end{tikzcd}

\note{les quatre flèches verticales portent chacune la même glose, lue
« isom. si $X'/X$ rig. » ; la flèche $u$ est annotée « surj. si $X'$ \ill{} »,
et la flèche $u'$ « surj. si $X'$ \ill{} ». Le membre médian de la ligne
supérieure porte en outre, d'une autre encre, la mention
\uncertain{« algèbre »} et une parenthèse \ill{}}

\[
  \varinjlim_{n} \struck{D^{\bullet}(i_{n})} \ill{}
\]
\note{un bloc de trois lignes, au milieu du feuillet, est entièrement
raturé ; deux relations subsistent en marge : $w'v' = \mathrm{id}$ et
$v'w' = \mathrm{id}$}

\[
  \operatorname{ch}\bigl(i_{!}\,x\bigr)\,\operatorname{Todd}(X'/S)
  \qquad X'/S
\]

\emph{Filtration sur $K^{\bullet}_{X}(X')$.}
\note{le feuillet s'arrête sur ce titre : rien ne suit}

\section*{R.R. des surfaces, et autres cas particuliers divers}

\note{titre pris du feuillet 39, qui ne porte que ces mots, de sa main : « R.R.
des \uncertain{surfaces}, et autres cas particuliers divers »}

\page{40}

$f : X \longrightarrow Y$ morphisme propre.

\begin{tikzcd}[column sep=large, row sep=small]
K(X) \arrow[r, "f_{!}"] \arrow[d, "\widetilde{c}_{X}"'] & K(Y) \arrow[d, "\widetilde{c}_{Y}"] \\
\widetilde{A}(X) \arrow[d, "h_{X}"'] & \widetilde{A}(Y) \arrow[d, "h_{Y}"] \\
A(X) \otimes \mathbb{Q} \arrow[d, "\tau_{h}(X)"'] & A(Y) \otimes \mathbb{Q} \arrow[d, "\tau_{h}(Y)"] \\
A(X) \otimes \mathbb{Q} \arrow[r, "f_{*}"] & A(Y) \otimes \mathbb{Q}
\end{tikzcd}

\[
  \tau_{h}(Y)\,\varphi_{Y}\bigl(f_{!}(\xi)\bigr)
  \;=\; f_{*}\bigl(\tau_{h}(X)\,\varphi_{X}(\xi)\bigr)
  \tag{1}
\]

\emph{Exemples.}

1) $f$ une injection. \ill{} $Y$ \ill{} alors \ill{} donnant
\[
  c^{d}_{Y}\bigl(f_{!}(\xi)\bigr) \;=\; f_{*}\bigl[\,c^{\ill{}}(\ill{})\,\bigr]
\]
\ill{} et (1) pour le \ill{} alors une formule meilleure,
\[
  \ill{}
  \;=\; f_{*}\left[\,
  \frac{c^{q+d}\bigl(\struck{\ill{}}\ \widetilde{C}(\xi)
    \ast \lambda_{-1}\bigl(\widetilde{C}(T_{Y/X})\bigr)\bigr)}
  {(-1)^{q}\,c^{q}_{Y}(T_{Y/X})}\,\right]
  \tag{2}
\]
\[
  \varphi_{Y}\bigl(f_{!}(\xi)\bigr)
  \;=\; f_{*}\Bigl(\varphi_{X}(\xi)\,\tau_{h}(T_{Y/X})^{-1}\Bigr)
  \tag{1 bis}
\]
\note{dans (2), la ligne inférieure est d'une encre bleue, écrite sous une
barre : elle est lue ici comme le dénominateur d'une fraction, mais elle peut
aussi bien être une variante portée après coup}

\add{(}5.13\add{)} \ill{} \add{(}$\leqslant$ Notes \uncertain{précises}\add{)}
\[
  f_{!}\bigl(\tau_{X}(F)\bigr) \;=\; \tau_{Y}(F) .
\]

2) $f$ \ill{} une \uncertain{surjection} fibrée. \ill{} $X$ \ill{} alors
\ill{}
\[
  f_{!}\bigl(\tau_{X}(E)\bigr) \;=\; \sum_{i} (-1)^{i}\,\tau_{Y}\bigl(E^{(i)}\bigr) ,
\]
\ill{} $E^{(i)}$ est le fibré vectoriel sur $Y$ dont \ill{} le $i$-ème
\ill{} de \ill{} du foncteur de $X \to Y$. On a donc
\[
  \varphi_{Y}\bigl(f_{!}(\tau_{X}(E))\bigr)
  \;=\; \sum_{i} (-1)^{i}\,\varphi_{Y}\bigl(E^{(i)}\bigr) .
\]
Dans le cas où $E^{(i)} = 0$ pour $i > 0$, \ill{} plat. \ill{} $E(u)$ \ill{}
provenant d'un \ill{} si \ill{} remplace $E$ \ill{} coh. fermé dans une
\ill{} \ill{} de $X$ \ill{} obtient ainsi $\varphi_{Y}\bigl(E^{(i)}\bigr)$
dans $\widetilde{C}_{Y}\bigl(E^{(i)}\bigr)$ \ill{} projectif\add{)} \ill{}
\ill{} 2 fortiori \ill{}
\note{les cinq dernières lignes sont traversées par l'encre du verso et ne se
laissent lire que par fragments}

\end{document}
